%=============================================================================
\section{The Dilution Ladder}
\label{sec:ladder}

The $k$NN estimator with $k_{\max}$ neighbors reaches scales up to the $k_{\max}$-th nearest-neighbor distance, which for the undiluted sample is a few times the mean inter-particle separation $\bar{d} = \nbar^{-1/3}$.  To extend the measurement to larger scales---where the correlation function is weak but cosmologically important (BAO, large-scale bias, etc.)---we use a geometric ladder of diluted subsamples.

\subsection{Construction}
\label{sec:ladder_construction}

Starting from the full data catalog of $N_d$ points, we partition them into a sequence of volume-filling subsamples.  At each dilution level $\ell$, the data are randomly divided (without replacement) into $R_\ell$ disjoint subsets, each containing $N_d / R_\ell$ points distributed throughout the full survey volume:
\begin{equation}
  R_0 = 1,\quad R_1 = 8,\quad R_2 = 64,\quad R_3 = 512,\quad \ldots\quad R_\ell = 8^\ell\,.
  \label{eq:dilution_levels}
\end{equation}
The dilution factor of 8 is chosen so that the mean inter-particle separation increases by a factor of $8^{1/3} = 2$ at each level.  Each subset requires its own tree, built on its $N_d/R_\ell$ points.  Because each subset spans the full survey volume, it can measure $\xi(r)$ at all separations up to the survey scale---the dilution increases the mean separation, bringing larger scales within reach of small $k$.

The $k$-th nearest neighbor at dilution $R_\ell$ probes a characteristic scale
\begin{equation}
  r_{\rm char}(k, R_\ell) \;\sim\; \Bigl(\frac{k\,R_\ell}{\nbar}\Bigr)^{1/3}\,.
  \label{eq:r_char}
\end{equation}

With $k = 1, \ldots, 8$ at each dilution level, the range of probed scales is:
\begin{equation}
  r_{\rm char}(1, R_\ell) \;\sim\; R_\ell^{1/3}\,\bar{d}\,,\quad r_{\rm char}(8, R_\ell) \;\sim\; 2\,R_\ell^{1/3}\,\bar{d}\,.
  \label{eq:scale_range}
\end{equation}
Each rung covers a factor of 2 in scale.  The factor-of-8 dilution between levels ensures that the $k = 1$ scale at level $\ell + 1$ matches the $k = 8$ scale at level $\ell$:
\begin{equation}
  r_{\rm char}(1, R_{\ell+1}) = r_{\rm char}(1, 8 R_\ell) = 2\,R_\ell^{1/3}\,\bar{d} = r_{\rm char}(8, R_\ell)\,.
  \label{eq:overlap}
\end{equation}
This produces complete overlap between adjacent rungs.

\subsection{Overlap as a consistency check}
\label{sec:overlap}

In the overlap region, two independent measurements of $\xi(r)$ are available:
\begin{itemize}[nosep]
  \item From level $\ell$, using $k = 8$: the 8th nearest neighbor in a sample of $N_d / R_\ell$ points.
  \item From level $\ell + 1$, using $k = 1$: the 1st nearest neighbor in a sample of $N_d / R_{\ell+1} = N_d / (8 R_\ell)$ points.
\end{itemize}
These two measurements use different subsets of the data, with different neighbor ranks and different survival corrections.  Agreement between them confirms that:
\begin{enumerate}[nosep]
  \item The $k_{\max} = 8$ truncation is adequate at level $\ell$ (no missing pairs).
  \item The survival correction (Section~\ref{sec:k1}) is handled correctly.
  \item The dilution procedure does not introduce biases (e.g.\ from edge effects or the interaction between dilution and the survey geometry).
\end{enumerate}
Disagreement in the overlap region is a red flag that triggers investigation before the measurement is trusted at larger scales.

\subsection{Variance estimation from the partition}
\label{sec:variance}

At each dilution level $R_\ell$, the $R_\ell$ disjoint subsamples are independent (no shared data points).  For each subsample $m$, the $k$NN estimator yields $\xihat_m(r)$.  The scatter across subsamples provides an empirical estimate of the variance:
\begin{equation}
  \widehat{\mathrm{Var}}[\xihat(r)] \;=\; \frac{1}{R_\ell - 1}\sum_{m=1}^{R_\ell} \bigl(\xihat_m(r) - \bar{\xi}(r)\bigr)^2\,,
  \label{eq:var_estimate}
\end{equation}
where $\bar{\xi}(r)$ is the mean over subsamples.  This is a direct, model-free estimate of the measurement variance, incorporating both Poisson noise and cosmic variance.

The degrees of freedom at each level are naturally well matched to the measurement.  At the lowest dilution (level 0, full data), $k = 1, \ldots, 8$ probes only small separations $r \lesssim 2\bar{d}$, where the number of independent spatial volumes in the survey is $V_{\rm survey}/r^3 \gg 1$.  No variance estimate from multiple subsamples is needed at this level---the measurement itself has negligible sample variance.  At the highest dilution (level $L$, $R_L$ subsamples), $k = 1, \ldots, 8$ probes scales $r \sim 2^L \bar{d}$, where cosmic variance is important---but there are $R_L$ independent measurements with $R_L - 1$ degrees of freedom.  The scale--dilution matching ensures that the variance estimate is well determined everywhere without special treatment.

\subsection{The full ladder}
\label{sec:full_ladder}

The complete measurement procedure is:

\begin{enumerate}[nosep]
  \item \textbf{Level 0} ($R_0 = 1$): use the full data catalog.  Measure DD and DR $k$NN-CDFs for $k = 1, \ldots, 8$.  Compute $\xihat(r)$ via Eq.~\eqref{eq:xi_estimator} from the mean separation $\bar{d}$ to $\sim 2\bar{d}$.  No variance estimate at this level (only one realization).
  \item \textbf{Level 1} ($R_1 = 8$): draw $M_1 \sim 8$ subsamples of $N_d / 8$ points.  For each, measure DD and DR $k$NN-CDFs for $k = 1, \ldots, 8$.  Compute $\xihat_m(r)$ for each subsample, covering $\sim 2\bar{d}$ to $\sim 4\bar{d}$.  Check overlap with Level 0 at $r \sim 2\bar{d}$.  Estimate variance from scatter across subsamples.
  \item \textbf{Level $\ell$} ($R_\ell = 8^\ell$): draw $M_\ell$ subsamples of $N_d / R_\ell$ points.  Measure, overlap-check, estimate variance.  Covers $\sim 2^\ell \bar{d}$ to $\sim 2^{\ell+1}\bar{d}$.
  \item Continue until the desired maximum scale is reached.
\end{enumerate}

For a survey with $N_d = 10^7$ and mean separation $\bar{d} \approx 5\,h^{-1}\mathrm{Mpc}$, four dilution levels reach $\sim 80\,h^{-1}\mathrm{Mpc}$, well past the BAO scale.  Five levels reach $\sim 160\,h^{-1}\mathrm{Mpc}$.

\subsection{Combining across levels}
\label{sec:combining}

The final $\xihat(r)$ is assembled from the ladder by selecting, at each $r$, the measurement from the level and $k$ value whose characteristic scale is closest to $r$ and whose CDF is near 0.5 (optimal sampling; Section~\ref{sec:errors}).  In the overlap regions, a weighted average of the two adjacent levels is used, with weights inversely proportional to the estimated variance.

The resulting $\xihat(r)$ is a continuous function of $r$, spanning from the smallest inter-particle scale to the largest scale accessible by the survey geometry, with an empirical error bar at every point.
